\section{Purpose}

\textbf{[STIPULATION]} FFGFT describes fundamental physics on a compact 4-torus
$T^4$ with the single parameter $\xi=4/30000$. The observable world is
$\mathbb{R}^3\times\mathbb{R}_t$. This document addresses the question of how
both are connected -- and shows that the chain $T^4\to T^0$ consists of
\emph{two qualitatively different} operations that must not be confused: a
\emph{duality decompactification} (nearly reversible, no mode losses) and a
\emph{geometric projection} (lossy, not reversible). A third operation, going in
the opposite direction, is added: the \emph{representational translation} to the
Hilbert space (bijective, lossless).

The Standard Model is situated therein as a \emph{decompactified projection}:
its constants ($\alpha$, $v$, Yukawa) are recursion outputs from $\xi$; whoever
takes them as inputs works in the unrolled picture, without naming the compact
whole. This is not a refutation of the Standard Model -- it computes correctly
in its domain; it is a location.

\section{The Tension in the Question}

\textbf{[B]} Three observations about the transition $T^4\to\mathbb{R}^3\times\mathbb{R}_t$
appear in mutual contradiction:

\begin{itemize}[leftmargin=2em,itemsep=3pt,topsep=3pt]
  \item \textbf{Loss picture.} Every projection $D\to D-1$ is not injective,
    loses information. The chain is not reversible.
  \item \textbf{Compactness picture.} $T^4$ is fully compact; the observable
    spacetime is decompactified.
  \item \textbf{Upward picture.} There is already a \emph{bijective} translation
    of $T^4$ into a higher-dimensional space -- the Hilbert space (A070) --
    which is \emph{lossless}.
\end{itemize}

\textbf{[B]} These three pictures cannot be naively overlaid: if the first step
were pure loss, no decompactified time direction could arise from it (forgetting
does not generate a new dimension); and it would remain puzzling how an upward
movement can be lossless while the downward movement destroys information. The
resolution: there are not one but \emph{three} dimension-changing operations of
\emph{two kinds}.

\section{The Structure of $T^4$}

\textbf{[B]} The fundamental level of FFGFT is the compact 4-torus
\[
  T^4 = S^1_x\times S^1_y\times S^1_z\times S^1_m,
\]
with three \emph{spatial} circles and one \emph{mass circle} $S^1_m$. The
fourth direction $S^1_m$ is not a mass dimension \emph{beside} a separate time,
but the \emph{substitute time direction} itself. Mass and time are two
\emph{readings} of one and the same dimension; the fundamental relation
$\tilde T\cdot m=1$ is the translation dictionary:
\[
  \underbrace{S^1_m\;\text{(compact)}}_{\text{mass reading, discrete spectrum}}
  \;\;\xleftrightarrow{\;\;\tilde T\cdot m=1\;\;}\;\;
  \underbrace{\mathbb{R}_t\;\text{(open)}}_{\text{time reading, continuous flow}}.
\]

\textbf{[K]} \textbf{Kaluza--Klein reading.} A compact circle of radius $R$
carries a discrete mass tower $m_n\sim n/R$. The four directions carry two radius
classes: the three spatial circles ($R_{x,y,z}\sim R_\text{large}$) generate a
quasi-continuous mode tower with $\sim0$ gap -- this becomes space; the mass
circle ($R_m\sim\xi\,\ell_P$) has a huge gap and a discrete, heavy spectrum --
this becomes the particle mass spectrum. The fundamental relation $\tilde T\cdot m=1$
is thereby structurally the Kaluza--Klein relation for the mass dimension.

\section{Type I -- Duality Decompactification: Mass Becomes Time}

\textbf{[B]} The first, qualitatively distinguished operation is the \emph{unrolling}
of the mass circle via its universal cover:
\[
  S^1_m\;\xrightarrow{\;\text{cover}\;}\ \mathbb{R}_t,
  \qquad p:\mathbb{R}_t\to S^1_m,\;\;t\mapsto e^{2\pi i\,t/R_m},
\]
governed by $\tilde T\cdot m=1$. This operation is \emph{embedding}, not
projection: it generates a new direction ($\mathbb{R}_t$) instead of discarding
one. Periodic structures on $S^1_m$ lift uniquely to $\mathbb{R}_t$; no mode
content is lost. The ``price'' of decompactification is not information loss but
the appearance of the time direction itself.

\textbf{[B]} What does not remain directly visible: the \emph{compactness}. In
the mass reading the topology is manifest -- the spectrum is discrete, the
identification $t\sim t+R_m$ directly expressed. In the time reading the same
compactness is only \emph{latent}: from continuous time alone the topological
identification is not uniquely reconstructible (the winding remains ambiguous).
Precisely there the recursion $\xi^n$ sets in -- it \emph{is} the spectral
signature by which the hidden compactness remains recognisable.

\textbf{[S]} \textbf{What makes unrolled time continuous and uniform.} Three
questions, three answers: (1) \emph{Why is $\mathbb{R}_t$ continuous?} Because
the unrolling is the universal cover $p:\mathbb{R}\to S^1$ -- a local
homeomorphism; the cover of a circle \emph{is} the continuous line. (2) \emph{Why
does time flow uniformly?} Because of the homogeneity of $S^1_m$ (translation
invariant, no distinguished point) and the single fixed eigenvalue $m=1/\tilde T$.
(3) \emph{Why is it unrolled at all?} Not the geometry decides this -- in the
compact $T^4$ all four circles are equivalent, there is no time. The
\emph{measurement process} (the embedded observer) chooses one circle as time and
unrolls it. Continuity is the forgotten winding: keeping it would give back the
discrete circle.

\section{Type II -- Geometric Projection: Forget Spatial Directions}

\textbf{[B]} All further steps ($D3\to D2\to D1\to D0$) are genuine geometric
projections of spatial directions -- a projection $\pi:\mathbb{R}^k\to\mathbb{R}^{k-1}$
with $\ker(\pi)=\mathbb{R}$. It collapses a complete fibre; the information about
the forgotten direction is not contained in the image. Without additional data
(holographic completeness, boundary conditions) not reconstructible. The recursion
$\xi^n$ does \emph{not} act on these steps -- spatial projection losses are genuine.

\textbf{[B]} The step $D3\to D2$ is mathematically the holographic principle;
in the horizon case the Bekenstein--Hawking bound $S\le A/(4\ell_P^2)$ appears here.

\section{Type III -- Representational Translation: Bijective Upward}

\textbf{[B]} Downward (geometric) and upward (representational) are
\emph{opposite} information properties:
\[
  T^4 \;\xleftrightarrow{\;\text{bijective}\;}\ \mathcal H = L^2(T^4)\otimes\mathbb{C}^3
\]
The translation (A070) is not an approximation but a bijective structural
statement: every QM computation is mappable into FFGFT language and vice versa.
\emph{No} physical extra dimensions are added -- the same 4D content is expressed
in a higher-dimensional \emph{language} in which it fits without reduction.
Losslessness follows precisely from the representational character.

\textbf{[B]} A geometric extension $T^4\to T^5\to\dots$ is excluded in FFGFT:
complexity arises from \emph{fractal depth} ($D_f=3-\xi$, self-similarity), not
from \emph{more} dimensions. This is the deliberate contrast to string theory
(10/11D, Calabi--Yau, landscape problem). ``Higher'' in FFGFT always means type
III (representational), never geometric.

\section{The Formal Chain, Level by Level}

\textbf{[B]} The chain $T^4\to T^0$ separated by type:

\begin{center}
\renewcommand{\arraystretch}{1.35}
{\small
\begin{tabular}{lllll}
\toprule
\textbf{Level} & \textbf{Operation} & \textbf{Type} & \textbf{Invariant} & \textbf{$\xi^n$} \\
\midrule
$D4\to D3{+}t$ & $S^1_m\to\mathbb{R}_t$ (unroll) & I & $\tilde Tm=1$ & spectral signature \\
$D3\to D2$ & forget space & II & area $A$ & --- \\
$D2\to D1$ & forget space & II & length $L$ & --- \\
$D1\to D0$ & forget space & II & $n\in\mathbb{Z}$ & --- \\
\bottomrule
\end{tabular}
}
\renewcommand{\arraystretch}{1.0}
\end{center}

\[
\underbrace{T^4\;\longrightarrow\;\mathbb{R}^3\times\mathbb{R}_t}_{\text{Type I: embedding, nearly reversible}}
\;\longrightarrow\;
\underbrace{\mathbb{R}^2\;\longrightarrow\;\mathbb{R}\;\longrightarrow\;\{\mathrm{pt}\}}_{\text{Type II: projecting, not reversible}}
\]

\section{Three Operations at a Glance}

\textbf{[B]}
\begin{center}
\renewcommand{\arraystretch}{1.35}
{\small
\begin{tabular}{llll}
\toprule
\textbf{Type} & \textbf{Operation} & \textbf{Character} & \textbf{Information} \\
\midrule
I & mass$\,\leftrightarrow\,$time ($\tilde Tm=1$) & geometric & nearly reversible \\
II & spatial projection (down) & geometric & lossy \\
III & $T^4\leftrightarrow\mathcal H$/Hilbert space (up) & representational & \textbf{bijective} \\
\bottomrule
\end{tabular}
}
\renewcommand{\arraystretch}{1.0}
\end{center}

The bijectivity of type III is the mirror image of the non-reversibility of type
II: geometrically one loses when reducing, representationally one loses nothing
when translating.

\section{The Role of the Recursion $\xi^n$}

\textbf{[K]} The recursion $\xi^n$ is \emph{not} a universal compensation for
projection losses -- its function is narrowly limited. For type I, $\xi^n$
generates the discrete mass spectrum (particle ladder) -- that is the compact mode
content of $S^1_m$, dual as a time-frequency structure. For type II, $\xi^n$ does
not act -- spatial projection losses are genuine and not touched by the recursion.

\section{Reversibility, Precisely Per Type}

\textbf{[B]} The type-I step is invertible in mode content as a duality
(Fourier/KK); the winding remains ambiguous -- ``nearly reversible''. The spatial
sub-chain $\mathbb{R}^3\to\{\mathrm{pt}\}$ is not invertible: there is no
continuous map $\rho:\{\mathrm{pt}\}\to\mathbb{R}^3$ with $\rho\circ\Pi=\mathrm{id}$
-- every $\pi_k$ has $\ker=\mathbb{R}\neq\{0\}$, the collapsed fibre structure is
no longer present in the image.

\section{Why Exactly Four Circles}

\textbf{[B]} No theory derives the dimension number of spacetime; FFGFT does
\emph{not derive the 4 forward} -- it takes $3+1$ as empirical input. The
question is therefore not ``why is the world four-dimensional?'' but: \emph{given
$3+1$ -- what is the minimal structure that carries everything?}

\textbf{[K]} Naively one would need \emph{five} building blocks: three for space,
one for a mass spectrum, one for time evolution. But the fundamental relation
$\tilde T\cdot m=1$ identifies mass reading and time reading of \emph{one}
dimension (type I): the fourth circle $S^1_m$ carries as a KK tower the mass
spectrum \emph{and}, decompactified, time. The duality lets one circle do the
work of two -- the naive $5$ collapses to $\mathbf{4}$. Fewer than four ($T^3$
alone): no mass spectrum, no time. More than four (a second dual circle): a second
independent mass tower -- superfluous, without correlate. ``Optimal'' means here
\emph{minimally sufficient}: the smallest dimension count in which space, full
mass spectrum, and time together arise from one geometric object.

\section{In the Compact Picture: the Time Circle Is Chosen by the Measurement}

\textbf{[S]} In the compact $T^4$ all four circles are identical: dimensionless
cyclic coordinates, no signature, no space/time distinction. The asymmetry (one
time, signature $(3,1)$) appears only under decompactification (type I): exactly
one circle unrolls via its universal cover into non-compact time, the other three
remain compact (discrete spectrum) or project (space). \emph{Which} circle becomes
time is no intrinsic geometric fact -- the geometry is symmetric. Time is
\emph{relational}: the selection is the measurement process itself. The embedded
observer reading from inside experiences one direction as duration; this
experiencing constitutes the time circle. The symmetry breaking is the measurement
process, not the geometry.

\section{The Topological Identification $t\sim t+R_m$ as Physical Observable}

\textbf{[K]} The topological identification is not a mere accompanying structure
but a physical observable. In the time reading its period is the fundamental
revolution time
\[
  T_\text{rev} = \frac{R_m}{c} = \frac{\xi\,\ell_P}{c}\approx 7{.}2\times10^{-48}\,\text{s},
  \qquad \Delta\phi = 2\pi\xi\approx 8{.}4\times10^{-4}\;\text{rad/revolution}.
\]

\textbf{[K]} Two testable readings: (1) \textbf{Mass reading (mature,
quantitative).} The discrete mode spectrum generates particle masses as torus
overtones; from this follow Koide, $g-2$, and three generations (A100/A110/A120).
Compactness is numerically redeemed. (2) \textbf{Time reading (falsifiable).}
The periodicity generates a CHSH deviation from the Tsirelson bound of order $\xi$;
for 73 qubits $\Delta\approx5{.}4\times10^{-4}$ (resolvable at $\sigma\sim10^{-4}$).
The IBM Heron r2 hardware test (May 2026) confirms Bell violation ($S=2{.}74$); the
$\xi$ effect is below the NISQ noise floor and not yet resolvable (A160). No
unfalsifiable accompanying structure.

\textbf{[S]} \textbf{Scaling: linear vs.\ logarithmic.} A naive \emph{coherent
linear accumulation} $r\cdot\Delta\phi$ over $r$ circuits would give physically
absurd phase drifts ($\sim10^{44}$ rad/s) and would contradict functioning atomic
clocks and deep circuits. The laboratory prediction scales \emph{logarithmically},
\[
  \delta\phi\sim\frac{\xi\,\ln(r)}{D_f}\approx2\times10^{-4}
  \quad (D_f=3-\xi),
\]
not linearly. The linear form is the worst-case bound of the RSA argument
(coherence over $r\approx N$ cycles), not the laboratory prediction.

\section{The Standard Model as Decompactified Projection}

\textbf{[K]} Every computation using $\alpha$, $v$, or Yukawa couplings carries
the $K_\text{frak}$ recursion \emph{indirectly} from the FFGFT perspective:
$\alpha=\xi\,E_0^2$ (A130), where $E_0$ already absorbs the $K_\text{frak}$ factor.
The masses are the $\xi^p$-ladder ($m=r\,\xi^p\,v$, A100); a Yukawa coupling
$g=m/v$ is already a $\xi$-ladder ratio. The Standard Model takes these quantities
as measured inputs; FFGFT \emph{derives} them. Working with $\alpha$, $v$, or
Yukawa in continuous fields on $\mathbb{R}$ is structurally working in the
\emph{unrolled} picture -- the type-I region of decompactification. The compact
origin of the constants is carried along, but not named.

\textbf{[S]} \textbf{What is established, what remains programme.} Established is
the containment relation for the sector FFGFT actually derives: $\alpha$, the
lepton masses, the Koide scalar $Q=2/3$, the phase $\theta=2/9$ (A120/A130) -- for
these it is shown that they are recursion outputs. Programme remains the \emph{full}
Yukawa matrix: quark masses, mixing angles (CKM/PMNS), CP phase. This separation
is the same as throughout the corpus: core derivations (proved from $\xi$), bridges
(algebraically proved), reductions (plausibility sketches).

\section{Open Questions}

\textbf{First, the large-scale outer scale is treated here as a declared external
input, not derived from $\xi$.} The two-type structure depends only on the
radius-class separation large/small, not on the numerical value. The connection
of the large-scale outer scale to $\ell_1$ remains open.

\textbf{Second, holographic completeness for type II ($D3\to D2$) is not worked
out.} To what extent can the type-II loss be reconstructed exactly (not just
partially) via boundary data?

\textbf{Third, the full Yukawa matrix remains programme.} The SM-subset relation
is a core derivation for the lepton sector; for quark masses/CKM/PMNS/CP-phase
it is an open obligation.

\section{Verification Scripts}

\begin{itemize}[leftmargin=2em,itemsep=2pt,topsep=2pt]
  \item Type-I reversibility: Fourier bijection $S^1_m\leftrightarrow\mathbb{R}_t$
    numerically, and winding ambiguity:
    \texttt{python/A\_Serie\_Skripte/a090\_projektion.py}
  \item CHSH deviation: $\Delta\text{CHSH}\approx\xi/(2\pi)$ numerically and
    logarithmic vs.\ linear regime scaling:
    \texttt{python/A\_Serie\_Skripte/a090\_chsh.py}
\end{itemize}

\section{Supersedes}

This document subsumes: Dok.~270 (projection chain $T^4\to T^0$, type I/II/III,
reversibility, role of $\xi^n$, why four circles, time circle as measurement
process, $t\sim t+R_m$ as observable, linear/logarithmic scaling), Dok.~298 (SM
as decompactified projection, Copernicus analogy, what is established/programme).
\emph{Not} taken over from Dok.~263 (fractal holography): the bulk with 31
cosmology hits, notably $H_0$ as projection artefact, vacuum energy, UIFT bridge
-- all do not belong in the A-series core (A010). Likewise not taken over from
Dok.~298: the HLV-YUKAWA section (IPI framework). The holographic projection
$D3\to D2$ (Dok.~263/268/254) is named but not carried out (open point).

\section{Use of AI Tools}

This document was created with the assistance of AI tools.
Responsibility for content and errors rests entirely with the author.
Full wording of the rules in A010.
